\section{Введение}
\label{sec:Chapter0} \index{Chapter0}

Современные программные системы нуждаются в больше количестве вычислений. Особенно остро данный вопрос стоит для серверных приложений, обрабатывающих значительные объёмы данных. Повышение производительности программного обеспечения позволяет эффективнее использовать аппаратные ресурсы, снижая потребность в дополнительном серверном оборудовании. Даже относительно небольшое ускорение может приводить к заметному экономическому эффекту в масштабах крупных систем.
Одним из важных аппаратных механизмов повышения производительности современных процессоров является SIMD-параллелизм (Single Instruction Multiple Data), позволяющий выполнять за одну инструкцию операцию одновременно над несколькими элементами данных. 
Несмотря на широкую аппаратную поддержку SIMD-инструкций в современных процессорных архитектурах, ручное написание векторизованного кода требует глубокого знания особенностей конкретной архитектуры,  что усложняет сопровождение программ и существенно снижает их переносимость. В связи с этим особую практическую значимость имеют методы автоматической векторизации, позволяющие компилятору самостоятельно преобразовывать скалярный код в его векторный эквивалент.
Особый интерес задача автоматической векторизации представляет для языка программирования Go. Данный язык широко применяется при разработке высоконагруженных серверных систем. Однако автоматическая векторизация в главной реализации компилятора Go отсутствует. Это приводит к недоиспользованию аппаратных возможностей современных процессоров и снижает потенциальную производительность программ.
Среди существующих подходов к автоматической векторизации важное место занимает метод Superword-Level Parallelism (SLP). В отличие от цикловой векторизации, ориентированной на анализ и преобразование циклических конструкций, SLP-векторизация направлена на поиск независимых однотипных операций внутри ациклических участков программы. Такой подход требует менее сложного анализа зависимостей, проще интегрируется в компилятор и обладает меньшими накладными расходами на этапе компиляции.
Выбор метода SLP особенно актуален для компилятора Go, так как реализация сложных цикловых оптимизаций потребовала бы существенного расширения существующей инфраструктуры анализа, тогда как SLP-подход позволяет достичь практического ускорения на ряде классов задач при умеренной сложности реализации. Сложность алгоритма особенно важна, поскольку одной из ключевых характеристик для разработчиков на Go является время компиляции.

